\markboth{EVALUATION}{EVALUATION}
\section{Evaluation Framework}\label{sec:evaluation}

Evaluating a multi-agent conversational system requires assessing several dimensions at once: does it call the right tools, ground answers in fetched data, maintain a supportive tone, and let the user achieve their goal across multiple turns? No single method covers all of this for an LLM-based system: a deterministic test suite cannot judge conversational quality, and a brief usability test cannot establish architectural reliability. The evaluation therefore combines four levels. \textit{Component and API stability tests} (Section~\ref{sec:evaluation:stability}) establish, independently of content quality, that the architecture itself works reliably; \textit{robustness tests} (Section~\ref{sec:evaluation:robustness}) actively probe up to which point the system keeps working; \textit{automated conversation-quality evaluation} (Sections~\ref{sec:evaluation:framework}--\ref{sec:evaluation:metrics}) scores the conversations real users had with the deployed system, using LLM judges and a deterministic grounding scorer; and a \textit{usability study} (Section~\ref{sec:evaluation:usability}) adds structured tasks and standardised usability measurement with the same users. One deliberate scope decision frames all four levels: assessing the sports-science \textit{quality} of a generated training plan is not a primary focus of this evaluation, as validating it would require sports-science expertise and longitudinal observation beyond the scope of this seminar project. The core evaluation object is instead whether the agentic system reliably selects and calls the correct MCP services and grounds its answers in the retrieved data.

\subsection{Component and API Stability Tests}\label{sec:evaluation:stability}

Before conversational quality can be assessed, the architecture beneath it must be shown to be stable: how often do tool calls fail, and does every call produce a response at all? The component-test suite answers this independently of any LLM. It exercises each of the eight native MCP servers directly through \texttt{ToolHost}, bypassing the agent layer entirely: every tool is invoked with representative arguments, and the suite records whether a well-formed MCP response arrives, whether its payload is valid JSON, and whether that payload is a domain result or a structured \texttt{\{"error": ...\}} object. Repeating the sweep against the live upstream APIs yields per-server response rates, success rates, and latency distributions (Table~\ref{tab:stability}, Appendix~\ref{sec:appendix:evaltables}). The suite is being consolidated from the development-time test scripts in the repository's \texttt{tests/} folder into a single reproducible run that doubles as a regression harness; its results are reported separately from the conversational scorers precisely because they isolate \textit{architectural} reliability from \textit{content} quality.

\subsection{Robustness and Failure Testing}\label{sec:evaluation:robustness}

Stability tests show the system works under normal conditions; robustness tests ask up to which point it keeps working, documenting not only that the system functions but where it fails and why. Each scenario in Table~\ref{tab:robustness} (Appendix~\ref{sec:appendix:evaltables}) deliberately disables or degrades one dependency (an MCP server, a specialist process, an upstream API, the LLM gateway) or stresses one boundary (malformed input, long conversations) and records the observed behaviour against the graceful-degradation design of Section~\ref{sec:agents:observability}. The distinction matters because the failure modes of an agentic system are not binary: between ``full answer'' and ``no answer'' lies a spectrum of degraded answers, and the tests document which degradations occur, whether the system reports them honestly to the user, and at which point it stops producing useful output altogether.

\subsection{Evaluation Pipeline}\label{sec:evaluation:framework}

The conversation-quality evaluation scores the conversations real people had with the deployed Training Copilot. An earlier iteration of this evaluation simulated synthetic user personas with an LLM role-playing the user; we replaced it with real-user logs because a persona simulator ultimately measures the system against an LLM's \textit{imitation} of a user, whereas logged conversations measure it against genuine user behaviour, vocabulary, and impatience -- at the cost of scripted control over domain coverage, which we account for in the reporting (Section~\ref{sec:evaluation:metrics}). We build on MLflow's multi-turn evaluation tooling~\citep{mlflow_multiturn_2025}: every chat turn is traced live into the user's own MLflow experiment and tagged with a session identifier, so turns group into complete conversations that session-level scorers can assess as a whole.

The LLM-judge scorers rest on the finding of \cite{zheng_judging_2023} that a strong model grading open-ended dialogue agrees with human preferences about as often as humans agree with each other, making LLM-as-a-judge a scalable proxy for human evaluation; their catalogue of judge biases (position, verbosity, self-enhancement) also motivates our model-separation safeguard (Section~\ref{sec:evaluation:models}). The design further draws on \cite{shankar_validates_2024}, who show LLM evaluators must be aligned with human preferences to be reliable, and our deterministic grounding scorer follows the $\tau$-bench approach of \cite{yao_taubench_2024}, which judges tool-agent interaction by comparing system state against annotated goals rather than chat text alone. The pipeline has three components: a \textbf{per-user trace store} (each user's conversations, recorded live during normal use); a set of \textbf{scorers} (an LLM judge for the conversational dimensions and a deterministic code checker for grounding); and an \textbf{experiment tracker} (MLflow) recording every conversation, trace, verdict, and metric. The pipeline is fully automated: a single command (\texttt{python -m evaluation.run\_users}) reads every user's logged conversations, scores them, and generates an HTML report with per-user scorecards.

\subsection{Real-User Conversation Corpus}\label{sec:evaluation:corpus}

The evaluation corpus is the set of conversations real users had with the deployed instance during the evaluation period. Participants were recruited from the authors' athletic circles, including athletically active testers from outside the project team, and were informed that their conversations are logged and scored for this evaluation. They used the system freely -- through the web app and the Telegram channel -- rather than following scripts, so the corpus reflects genuine information needs across the specialist domains: recovery checks, training-load questions, route planning, weather decisions, and training-plan requests. Who these users are is captured by the usability study's intake questionnaire (Section~\ref{sec:evaluation:usability}); the aggregated participant statistics (demographics, sporting background, fitness-app and AI usage) and the corpus usage statistics (conversations, messages, turns per conversation, channel split) are reported as charts in Appendix~\ref{sec:appendix:ux} (Figures~\ref{fig:ux_participants} and~\ref{fig:corpus_usage}), pending the end of the logging period.

Because usage is free-form, domain coverage is not guaranteed by construction the way scripted test cases would guarantee it; we therefore report which specialist domains the logged conversations actually exercised as part of the results, rather than assuming full coverage.

\subsection{Scoring Dimensions}\label{sec:evaluation:scorers}

Each conversation is scored along five dimensions, combining an LLM judge with a deterministic code scorer (Table~\ref{tab:scorers}, Appendix~\ref{sec:appendix:evaltables}).

The fifth scorer, \texttt{grounded\_in\_real\_data}, inspects the tool-call \textit{spans} reconstructed in each turn's MLflow trace rather than asking an LLM to guess whether tools were used: it counts tool calls, identifies which were invoked and whether they succeeded, and produces a per-turn breakdown. The verdict is binary (``yes'' if at least one tool was called across the conversation), but the rationale gives granular visibility into grounding quality; a factual ``did it call tools or not?'' question is far more reliably answered from the trace than guessed from chat text. The \texttt{supportive\_coaching\_tone} dimension asserts that the assistant stay supportive even when the user is impatient, sceptical, or pushing for specifics, never dismissive or robotic. This tests resilience under the demanding user behaviour \cite{barger_artificial_2025} identified as important for AI-coaching acceptance -- and unlike a scripted simulation, real users supply that behaviour unprompted.

Because an LLM judge is only as reliable as its prompt, the judge's inputs and assertions are defined explicitly rather than left to a generic quality prompt: the judge receives the full transcript and rates the four conversational dimensions against fixed verdict definitions (completeness, frustration, safety, supportive tone), returning structured verdicts with a short rationale, while \texttt{grounded\_in\_real\_data} is deterministic code over the trace and never consults an LLM. A planned refinement calibrates the judge against annotated positive and negative example conversations drawn from the corpus itself, which then serve as scoring anchors, so that a verdict is anchored to explicit criteria rather than the judge's own latent preferences~\citep{zheng_judging_2023, shankar_validates_2024}.

One capability these five scorers do not yet isolate is \textit{citation faithfulness}: whether the fitness specialist's evidence-based claims are actually supported by the passages it retrieved, rather than fluently invented. Logged conversations in which users explicitly ask for sources exercise this qualitatively, but a dedicated deterministic scorer, checking that every citation in the answer matches a passage present in that turn's RAG-tool span (the way \texttt{grounded\_in\_real\_data} reads MCP spans), is a planned sixth dimension we report as future work. Alongside the automated scorers, the one-click \textit{report button} (Section~\ref{sec:agents:observability}) lets users flag a response and captures its trace references; turning these flags into a labelled good/bad validation set, against which the LLM judges themselves can be checked for the alignment \cite{shankar_validates_2024} argue is needed before trusting them, is future work.

\subsection{Trace Reconstruction}\label{sec:evaluation:trace}

A technical challenge in evaluating the multi-agent system is that the deep tool-call spans produced by the agent server processes live in separate MLflow experiments from the per-user conversation traces. The per-user tracking layer addresses this by reconstructing the Copilot's specialist and tool-call structure from the \texttt{trace} dict returned by \texttt{orchestrator.run()} and emitting it as real child spans within the turn's trace (Figure~\ref{fig:trace_reconstruction}): each chat turn becomes one trace in the user's own experiment, with a root span carrying the user message and the answer, an agent span per consulted specialist, and a tool span per MCP call.

\begin{figure}[ht]
\centering
\resizebox{\textwidth}{!}{%
\begin{tikzpicture}[
    node distance=0.3cm and 0.5cm,
    every node/.style={font=\small},
    spanbox/.style={draw, thick, rounded corners=2pt, minimum height=0.55cm, align=center, font=\small},
    rootspan/.style={spanbox, draw=tcblue, fill=tcblue!8, minimum width=9cm},
    agentspan/.style={spanbox, draw=tcyellow!80!black, fill=tcyellow!8, minimum width=4cm},
    toolspan/.style={spanbox, draw=tccyan, fill=tccyan!6, minimum width=3.2cm},
    arr/.style={-{Stealth[length=4pt]}, thick, tcgrey!50},
]
    \node[rootspan] (root) {\texttt{fitdash\_copilot} (AGENT span)};

    \node[agentspan, below left=0.6cm and 0.3cm of root] (ag1) {\texttt{agent:recovery}};
    \node[agentspan, below right=0.6cm and 0.3cm of root] (ag2) {\texttt{agent:context}};

    \node[toolspan, below left=0.5cm and -0.2cm of ag1] (t1) {\texttt{garmin\_\_sleep}};
    \node[toolspan, right=0.15cm of t1] (t2) {\texttt{garmin\_\_hrv}};

    \node[toolspan, below left=0.5cm and -0.2cm of ag2] (t3) {\texttt{weather\_\_forecast}};
    \node[toolspan, right=0.15cm of t3] (t4) {\texttt{calendar\_\_events}};

    \draw[arr] (root.south) -- (ag1.north);
    \draw[arr] (root.south) -- (ag2.north);
    \draw[arr] (ag1.south) -- (t1.north);
    \draw[arr] (ag1.south) -- (t2.north);
    \draw[arr] (ag2.south) -- (t3.north);
    \draw[arr] (ag2.south) -- (t4.north);

\end{tikzpicture}%
}% end resizebox
\caption[Reconstructed span tree in the evaluation trace]{Reconstructed span tree in the evaluation trace. Agent and tool spans are rebuilt from the Copilot's \texttt{trace} dict and nested under the root evaluation span. Each tool span carries four custom attributes: \texttt{fitdash.tool} (name), \texttt{fitdash.tool\_ok} (success), \texttt{fitdash.tool\_error} (message), and \texttt{fitdash.duration\_ms} (latency).}\label{fig:trace_reconstruction}
\end{figure}


Each reconstructed tool span carries four custom attributes: \texttt{fitdash.tool} (the namespaced tool name), \texttt{fitdash.tool\_ok} (boolean success), \texttt{fitdash.tool\_error} (error message if any), and \texttt{fitdash.duration\_ms} (tool-call latency). The grounding scorer reads only these, never the chat text, so its verdict is deterministic and model-independent.

\subsection{Model Separation}\label{sec:evaluation:models}

A critical design decision is separating the models used for evaluation from the model under test. The evaluation harness uses official OpenAI models: GPT-5.4-nano for the LLM judge and GPT-5.4-mini for report generation. Training Copilot itself also runs on official OpenAI; we tried the KIT-hosted gateways (the KI Toolbox and the DSI Experimente gateway) but found them too slow for interactive use. So the judge and the system under test share a provider but not a model: the judge is a small, dedicated model (GPT-5.4-nano) that differs from the one driving the system under test. The \texttt{apply\_openai\_routing()} function points the evaluation process at the official OpenAI key and base URL.

\subsection{Metrics and Reporting}\label{sec:evaluation:metrics}

Each judged dimension yields a per-conversation verdict, which we aggregate as a \textit{pass rate} (the fraction of logged conversations the judge marks acceptable), reported overall and per user; the deterministic \texttt{grounded\_in\_real\_data} scorer aggregates the same way over its binary per-conversation verdict. Alongside the five quality dimensions we log operational measures per conversation (the number of turns, turns with tool errors, and end-to-end latency) so conversational quality can be read against cost, and we report which specialist domains the corpus exercised (Section~\ref{sec:evaluation:corpus}). MLflow rolls these into experiment-level metrics, and the generated HTML report renders per-user scorecards next to every transcript and trace for inspection. Table~\ref{tab:results} (Appendix~\ref{sec:appendix:evaltables}) gives the reporting structure; its cells are placeholders pending the final run. We stress that a small real-user corpus is a \textit{formative} sample: it reflects genuine usage rather than scripted coverage, is sized to surface failure modes rather than to establish statistical significance, and we report it as a characterisation of system behaviour rather than a confirmatory benchmark.

\subsection{Usability Study}\label{sec:evaluation:usability}

The automated pipeline scores free-form conversations; a structured usability study complements it. Each session is framed by a questionnaire and three hands-on tasks. The questionnaire's pre-task blocks capture the participant profile: demographics, sporting background (training frequency, ambition, social context, competition experience, and whether -- and from where -- they follow a training plan), fitness-app and wearable usage, and prior AI exposure (usage frequency, self-rated familiarity, and trust). Participants then use Training Copilot to accomplish three predefined tasks that walk through the core coaching loop end to end: (i) add an upcoming race -- the Baden Marathon Karlsruhe -- to the app, (ii) have the coach generate a training plan preparing for it, and (iii) plan today's workout (route, timing, and consistency with that plan). Afterwards, they complete the ten-item System Usability Scale~\citep{brooke_sus_1996} and closing questions on the overall experience: what they liked, what they did not, suggested improvements, and missing features. The full questionnaire structure and the aggregated results are reported in Appendix~\ref{sec:appendix:ux} (Figures~\ref{fig:ux_participants}--\ref{fig:ux_sus}, placeholders pending the study's completion); Table~\ref{tab:usability} (Appendix~\ref{sec:appendix:evaltables}) shows the quantitative reporting structure. The same participants' free-form conversations with the deployed instance flow into the conversation corpus of Section~\ref{sec:evaluation:corpus}, so the automated scorers assess the same user group whose usability perceptions the SUS captures. Like the conversation corpus, this is a \textit{formative} study: it is sized to surface usability problems and ground the SUS score, not to establish statistical significance.

\subsection{Expected Outcomes}\label{sec:evaluation:expected}

We expect high grounding (the orchestrator has no MCP access of its own, so specialists must call tools) and coaching-tone resilience even under impatient or sceptical users; because real usage is free-form, we additionally expect uneven domain coverage and report the domains the corpus actually exercised alongside the scores. The likely weak spots are multi-step sequential queries (latency, occasional context loss) and the historical fitness corpus. Quantitative results follow in the final submission and will guide further tuning of the prompts and specialist scoping.
